\section{Introduction}
\label{sec:intro}

Retrieval-augmented generation~\cite{lewis2020rag} lets a language model answer from a document
collection it was not trained on. In enterprise deployments the collection is shared by several
tenants and roles, so the application decides per request which documents a user may see. The
common design copies document ACLs from the identity system into the vector index and filters
retrieval on them~\cite{jeong2025permissionrag,chen2025integrating,zhong2026honeybee,doc_azure_trimming}.

That design makes the identity system the source of truth and every other component a copy, and
a RAG application makes several copies that outlive the request that created them. The index holds
materialized ACLs, refreshed by a crawl or a webhook~\cite{doc_graph_connectors}. A semantic cache
returns stored answers to similar questions~\cite{bang2023gptcache,doc_redisvl}. Chat memory keeps
earlier turns and, after a few turns, a summary written by the model, which names no source
document~\cite{doc_langchain_summary,doc_llamaindex_memory,packer2023memgpt}. When an administrator
revokes access, the identity system changes at once; each copy changes only when something
propagates the revocation to it. We call what the copies serve in the meantime
\emph{post-revocation leakage}.

\paragraph{Why it matters if the user saw it before.}
A revoked user has, by definition, been allowed to read the document, so an obvious objection is
that re-showing it reveals nothing. We take the objection seriously and measure around it. First,
we separate \emph{re-disclosure} (secrets the user was shown before revocation) from \emph{new
disclosure} (secrets of the document the user was never shown). Only the second is new information; the first is still a failure of complete
mediation~\cite{saltzer1975protection}, which access-control standards require to be closed
promptly when a user changes role or leaves~\cite{nist80053r5}, and which matters for users who no
longer hold their copies. Second, derived state is shared: a cache entry or a laundered answer can
reach a colleague who never had access. Third, a user who changes department (a \emph{mover})
keeps an active account, so the window is open to an ordinary, authenticated client.

\paragraph{Existence inference.}
A separate question is whether the application reveals that a restricted document \emph{exists}
without returning its content. Relevance scores, reranker confidence, latency and refusal wording
can all depend on a hidden document~\cite{buttcher2005security,gelernter2015xs}, and caches and LLM
serving have their own timing channels~\cite{gu2025auditing,song2025earlybird,zheng2024inputsnatch,carlini2024remote}.
We measure existence inference as a distinguishing game between two worlds that differ only in the
target document, for insiders (same tenant) and outsiders, and keep it separate from leakage
throughout.

\paragraph{Research questions.}
\begin{itemize}
\item[RQ1] Which derived artifacts disclose a revoked document, and how much of what they disclose
is new to the user?
\item[RQ2] How long does each leak last in the day after revocation, under a daily ACL crawl and
default cache and session settings?
\item[RQ3] What can an insider or outsider infer from response metadata about documents they cannot
read?
\item[RQ4] Which mitigations close each channel? How do read-time defenses for derived state
(provenance taint, permission epochs, session flush) and published retrieval-time designs compare?
\item[RQ5] What do the defenses cost the revoked user on documents it still holds, and in
authorization load?
\end{itemize}

\paragraph{Contributions.}
\begin{enumerate}
\item A measurement procedure for post-revocation leakage per derived artifact, with paired
revocation events, right-censored lifetimes, read-only probes that do not refresh what they
measure, and a decomposition into re-disclosure and new disclosure (\cref{sec:method}).
\item Measurements on real e-mail (Enron~\cite{klimt2004enron,ryan2025enronqa}) with
Llama-3.1-8B (4-bit, served locally), whose chat memory is summarized by the same model, and a
check of the two main leaks in an unmodified open-source stack (\cref{sec:results}).
\item A twin-world measurement of existence inference that separates insiders from outsiders and
content channels from metadata channels, with intervals clustered by target
document~\cite{obuchowski1997clustered}, run with a deterministic stub generator (\cref{sec:rq3}).
\item A comparison of read-time defenses for the derived state of a RAG pipeline (provenance taint
and a per-user permission epoch) against a published retrieval-time design, on leakage, utility for
the revoked user, and authorization load, with session flush and further designs in the stub harness
(\cref{sec:defense,sec:rq4,sec:rq5}). The mechanisms are known
(information-flow control~\cite{myers1997decentralized,elnikety2016thoth,costa2026fides}, complete
mediation); the contribution is their instantiation for LLM-derived artifacts and the measured
trade-off between them.
\end{enumerate}

\paragraph{Summary of findings.}
\pending{Fill after the runs of configs/min/, at most five numbers: LM and LM\_new under baseline
(p\_min\_enron), the retrieval-time design's residual leakage, the retained-document recall of epoch
against full, insider existence AUC under baseline and under the defense (stub).}

\paragraph{Scope.}
The deployment model is grounded in documented framework and product behaviour
(\cref{tab:grounding}), but every measurement comes from our testbed rather than from a production
system, apart from the real-stack check (\cref{sec:res-realstack}). We use one 4-bit 8B model and
50 revocation events per cell. Zero measured leakage means zero under a corpus-unique exact-match
judge and an unvalidated model judge (\cref{sec:method}), for an adversary who queries the system.
\cref{sec:limits} states where that does not hold.
